% !TeX root = ../presentacion.tex
% ===========================================================================
%  8. Reportes — 0:35
%
%  Criterio: «Módulo de reportes básicos» (10 %). La rúbrica pide que los
%  datos sean «consistentes respecto a las reservas registradas»: por eso la
%  demo crea una reserva y vuelve a esta pantalla a ver el número moverse.
%
%  La captura sale de ../informe/figuras/capturas (las mismas del E1).
% ===========================================================================

\bloque{Reportes}{0:35}{\exponeFrontend}%
       {Módulo de reportes básicos (10\,\%)}

\begin{frame}{Los cuatro indicadores, sobre el rango que elija el administrador}
  \note{Enumerar rápido y anunciar la prueba de consistencia: en la demo se
  crea una reserva y se vuelve aquí a ver el número subir. Es lo que separa
  un reporte de una pantalla con datos de adorno.}

  \begin{columns}[T]
    \begin{column}{0.42\textwidth}
      \begin{enumerate}\scriptsize
        \item Reservas por cancha y por deporte.
        \item Ocupación por cancha: horas reservadas sobre horas disponibles.
        \item Cancelaciones del período.
        \item Canchas de mayor y menor demanda.
      \end{enumerate}
      \vspace{2pt}
      {\scriptsize Solo lectura, solo administrador. Sin almacén analítico:
      \id{ms-reportes} compone por REST lo que ya está en \id{ms-reservas} y
      \id{ms-canchas}.}\\[4pt]
      {\scriptsize\color{textoSuave} En la demo se crea una reserva y se
      vuelve aquí: \alert{el número se mueve}.}
    \end{column}
    \begin{column}{0.56\textwidth}
      \includegraphics[width=\linewidth]{capturas/admin-reportes-indicadores}
    \end{column}
  \end{columns}
\end{frame}
